% !TeX root = ../Trabajo_RPM.tex
\clearpage
\section{Inventario y resultados de las cien instancias}\label{an:instancias100}
Las tablas de este anexo incluyen los cien problemas registrados antes de resolver. Pub. identifica la matriz publicada y Rec. la reconstruida. N conserva la capacidad nominal; M1--25 indica la reducción del 25\,\% en esa máquina. Los cuatro identificadores consecutivos de cada celda comparten demanda, tamaño y horizonte. La columna de huella muestra los primeros doce caracteres del SHA-256 del archivo de entrada; el manifiesto digital contiene los 64 caracteres. Ningún problema se retiró por su resultado.

\subsection{Diseño e identidad de las entradas}
{\footnotesize\setlength{\tabcolsep}{4pt}
\begin{longtable}{lrrrlrrl}
\caption{Inventario de las cien instancias sintéticas}\label{tab:inventario100}\\
\toprule
ID & $n$ & $H$ & $\lambda$ & Cap. & Matriz & $\sum_iD_i$ & SHA-256 (prefijo) \\
\midrule\endfirsthead
\multicolumn{8}{l}{\tablename\ \thetable\ (continuación)}\\\toprule
ID & $n$ & $H$ & $\lambda$ & Cap. & Matriz & $\sum_iD_i$ & SHA-256 (prefijo) \\
\midrule\endhead\bottomrule\endfoot
E001 & 99 & 1 & 0,8 & N & Pub. & 10629 & \texttt{ff760c1c1ec2} \\
E002 & 99 & 1 & 0,8 & N & Rec. & 10629 & \texttt{534ded157223} \\
E003 & 99 & 1 & 0,8 & M1--25 & Pub. & 10629 & \texttt{c5b207fcd43f} \\
E004 & 99 & 1 & 0,8 & M1--25 & Rec. & 10629 & \texttt{5a5bdb9d54aa} \\
E005 & 99 & 3 & 0,8 & N & Pub. & 17027 & \texttt{75b4182b36df} \\
E006 & 99 & 3 & 0,8 & N & Rec. & 17027 & \texttt{794917d326fb} \\
E007 & 99 & 3 & 0,8 & M1--25 & Pub. & 17027 & \texttt{6b08808edeaa} \\
E008 & 99 & 3 & 0,8 & M1--25 & Rec. & 17027 & \texttt{594ff9cb27b6} \\
E009 & 150 & 1 & 0,8 & N & Pub. & 8844 & \texttt{8befae5b96c4} \\
E010 & 150 & 1 & 0,8 & N & Rec. & 8844 & \texttt{1504370e0275} \\
E011 & 150 & 1 & 0,8 & M1--25 & Pub. & 8844 & \texttt{0dd05916ab41} \\
E012 & 150 & 1 & 0,8 & M1--25 & Rec. & 8844 & \texttt{a3fe3d07ea19} \\
E013 & 150 & 5 & 0,8 & N & Pub. & 71235 & \texttt{f4b4f24415e7} \\
E014 & 150 & 5 & 0,8 & N & Rec. & 71235 & \texttt{cdc613095992} \\
E015 & 150 & 5 & 0,8 & M1--25 & Pub. & 71235 & \texttt{b5022fcdb63e} \\
E016 & 150 & 5 & 0,8 & M1--25 & Rec. & 71235 & \texttt{d09cbec711fe} \\
E017 & 200 & 3 & 0,8 & N & Pub. & 58103 & \texttt{afa8609450b2} \\
E018 & 200 & 3 & 0,8 & N & Rec. & 58103 & \texttt{b520b38f6d05} \\
E019 & 200 & 3 & 0,8 & M1--25 & Pub. & 58103 & \texttt{9f86b70f2107} \\
E020 & 200 & 3 & 0,8 & M1--25 & Rec. & 58103 & \texttt{c605376ef614} \\
E021 & 200 & 5 & 0,8 & N & Pub. & 101236 & \texttt{7fa35d29b51e} \\
E022 & 200 & 5 & 0,8 & N & Rec. & 101236 & \texttt{8a1e64bcc33f} \\
E023 & 200 & 5 & 0,8 & M1--25 & Pub. & 101236 & \texttt{6ac6c873bff6} \\
E024 & 200 & 5 & 0,8 & M1--25 & Rec. & 101236 & \texttt{e2d31d2553d6} \\
E025 & 99 & 3 & 1,0 & N & Pub. & 37073 & \texttt{2b3db351ecbe} \\
E026 & 99 & 3 & 1,0 & N & Rec. & 37073 & \texttt{3b2b141341cf} \\
E027 & 99 & 3 & 1,0 & M1--25 & Pub. & 37073 & \texttt{a46ffc432020} \\
E028 & 99 & 3 & 1,0 & M1--25 & Rec. & 37073 & \texttt{3936c6852a83} \\
E029 & 99 & 5 & 1,0 & N & Pub. & 54048 & \texttt{5a70ee50cc87} \\
E030 & 99 & 5 & 1,0 & N & Rec. & 54048 & \texttt{64c238c3581a} \\
E031 & 99 & 5 & 1,0 & M1--25 & Pub. & 54048 & \texttt{9dbeba79390a} \\
E032 & 99 & 5 & 1,0 & M1--25 & Rec. & 54048 & \texttt{03f71cefad69} \\
E033 & 150 & 1 & 1,0 & N & Pub. & 13261 & \texttt{4f9b628c4ef2} \\
E034 & 150 & 1 & 1,0 & N & Rec. & 13261 & \texttt{bfccd016470b} \\
E035 & 150 & 1 & 1,0 & M1--25 & Pub. & 13261 & \texttt{4c0bd7f18bb3} \\
E036 & 150 & 1 & 1,0 & M1--25 & Rec. & 13261 & \texttt{e734cb1b0758} \\
E037 & 150 & 3 & 1,0 & N & Pub. & 40919 & \texttt{d5ffedc77e88} \\
E038 & 150 & 3 & 1,0 & N & Rec. & 40919 & \texttt{6876d505eb79} \\
E039 & 150 & 3 & 1,0 & M1--25 & Pub. & 40919 & \texttt{0d49d712c721} \\
E040 & 150 & 3 & 1,0 & M1--25 & Rec. & 40919 & \texttt{430969ca97f3} \\
E041 & 200 & 1 & 1,0 & N & Pub. & 15286 & \texttt{e0231dcb8344} \\
E042 & 200 & 1 & 1,0 & N & Rec. & 15286 & \texttt{b63bcacd9a1b} \\
E043 & 200 & 1 & 1,0 & M1--25 & Pub. & 15286 & \texttt{044522bd3c69} \\
E044 & 200 & 1 & 1,0 & M1--25 & Rec. & 15286 & \texttt{5944a3d0c980} \\
E045 & 200 & 5 & 1,0 & N & Pub. & 80014 & \texttt{e50a9264137f} \\
E046 & 200 & 5 & 1,0 & N & Rec. & 80014 & \texttt{44bc8b155c14} \\
E047 & 200 & 5 & 1,0 & M1--25 & Pub. & 80014 & \texttt{0e7e347093de} \\
E048 & 200 & 5 & 1,0 & M1--25 & Rec. & 80014 & \texttt{7894220ef002} \\
E049 & 99 & 1 & 1,2 & N & Pub. & 6931 & \texttt{138251c387e3} \\
E050 & 99 & 1 & 1,2 & N & Rec. & 6931 & \texttt{b5506f205db7} \\
E051 & 99 & 1 & 1,2 & M1--25 & Pub. & 6931 & \texttt{56e9d62fe561} \\
E052 & 99 & 1 & 1,2 & M1--25 & Rec. & 6931 & \texttt{085cb3cc14e9} \\
E053 & 99 & 5 & 1,2 & N & Pub. & 29277 & \texttt{4f04e35d805c} \\
E054 & 99 & 5 & 1,2 & N & Rec. & 29277 & \texttt{2700bf324179} \\
E055 & 99 & 5 & 1,2 & M1--25 & Pub. & 29277 & \texttt{9bf39c3008f5} \\
E056 & 99 & 5 & 1,2 & M1--25 & Rec. & 29277 & \texttt{e7fae9406fc5} \\
E057 & 150 & 3 & 1,2 & N & Pub. & 26535 & \texttt{c4cf5a947936} \\
E058 & 150 & 3 & 1,2 & N & Rec. & 26535 & \texttt{db54dcecd27b} \\
E059 & 150 & 3 & 1,2 & M1--25 & Pub. & 26535 & \texttt{6c40e5738d7d} \\
E060 & 150 & 3 & 1,2 & M1--25 & Rec. & 26535 & \texttt{2f9e1a8e1e72} \\
E061 & 150 & 5 & 1,2 & N & Pub. & 64183 & \texttt{4667f6609a52} \\
E062 & 150 & 5 & 1,2 & N & Rec. & 64183 & \texttt{6c1ab7b5c707} \\
E063 & 150 & 5 & 1,2 & M1--25 & Pub. & 64183 & \texttt{9e4141e0a8f0} \\
E064 & 150 & 5 & 1,2 & M1--25 & Rec. & 64183 & \texttt{4c144aeef1c7} \\
E065 & 200 & 1 & 1,2 & N & Pub. & 28035 & \texttt{21c43665f934} \\
E066 & 200 & 1 & 1,2 & N & Rec. & 28035 & \texttt{e4a69936a901} \\
E067 & 200 & 1 & 1,2 & M1--25 & Pub. & 28035 & \texttt{0132978ca3a8} \\
E068 & 200 & 1 & 1,2 & M1--25 & Rec. & 28035 & \texttt{c95b49401ec5} \\
E069 & 200 & 3 & 1,2 & N & Pub. & 50673 & \texttt{580f922deb1f} \\
E070 & 200 & 3 & 1,2 & N & Rec. & 50673 & \texttt{bd5499124761} \\
E071 & 200 & 3 & 1,2 & M1--25 & Pub. & 50673 & \texttt{5be12956a533} \\
E072 & 200 & 3 & 1,2 & M1--25 & Rec. & 50673 & \texttt{73fc939dc8d0} \\
E073 & 99 & 1 & 1,4 & N & Pub. & 6951 & \texttt{fe729f00b602} \\
E074 & 99 & 1 & 1,4 & N & Rec. & 6951 & \texttt{b7b188647342} \\
E075 & 99 & 1 & 1,4 & M1--25 & Pub. & 6951 & \texttt{a5be0eeae3a5} \\
E076 & 99 & 1 & 1,4 & M1--25 & Rec. & 6951 & \texttt{0b73015429f2} \\
E077 & 99 & 5 & 1,4 & N & Pub. & 54472 & \texttt{e640b256e4fe} \\
E078 & 99 & 5 & 1,4 & N & Rec. & 54472 & \texttt{46226094dcd5} \\
E079 & 99 & 5 & 1,4 & M1--25 & Pub. & 54472 & \texttt{6c83b348d468} \\
E080 & 99 & 5 & 1,4 & M1--25 & Rec. & 54472 & \texttt{68f9092556bd} \\
E081 & 150 & 3 & 1,4 & N & Pub. & 49777 & \texttt{9ef0d86b4922} \\
E082 & 150 & 3 & 1,4 & N & Rec. & 49777 & \texttt{af6ca44d8006} \\
E083 & 150 & 3 & 1,4 & M1--25 & Pub. & 49777 & \texttt{9ef664273b28} \\
E084 & 150 & 3 & 1,4 & M1--25 & Rec. & 49777 & \texttt{90788c9c9f08} \\
E085 & 150 & 5 & 1,4 & N & Pub. & 79823 & \texttt{a87623c84f63} \\
E086 & 150 & 5 & 1,4 & N & Rec. & 79823 & \texttt{6085548c1618} \\
E087 & 150 & 5 & 1,4 & M1--25 & Pub. & 79823 & \texttt{30845f82a75e} \\
E088 & 150 & 5 & 1,4 & M1--25 & Rec. & 79823 & \texttt{b4a98dd7cb13} \\
E089 & 200 & 1 & 1,4 & N & Pub. & 16536 & \texttt{f625fe17aedc} \\
E090 & 200 & 1 & 1,4 & N & Rec. & 16536 & \texttt{28fddd643596} \\
E091 & 200 & 1 & 1,4 & M1--25 & Pub. & 16536 & \texttt{cd63851f98c0} \\
E092 & 200 & 1 & 1,4 & M1--25 & Rec. & 16536 & \texttt{4321945fa34d} \\
E093 & 200 & 3 & 1,4 & N & Pub. & 77174 & \texttt{ddb5472d7b3c} \\
E094 & 200 & 3 & 1,4 & N & Rec. & 77174 & \texttt{8f3336d3f495} \\
E095 & 200 & 3 & 1,4 & M1--25 & Pub. & 77174 & \texttt{93389625232c} \\
E096 & 200 & 3 & 1,4 & M1--25 & Rec. & 77174 & \texttt{fad24f18313c} \\
E097 & 200 & 5 & 1,4 & N & Pub. & 69892 & \texttt{168353935317} \\
E098 & 200 & 5 & 1,4 & N & Rec. & 69892 & \texttt{ec402ae149c0} \\
E099 & 200 & 5 & 1,4 & M1--25 & Pub. & 69892 & \texttt{5277f11544a9} \\
E100 & 200 & 5 & 1,4 & M1--25 & Rec. & 69892 & \texttt{fbe833e0793b} \\
\end{longtable}

}

\subsection{Estado de cada método}
Ópt. significa óptimo acreditado por estado, cotas y verificación; F, plan factible sin cierre; Inv., inviabilidad declarada de forma concordante; y SI, ausencia de incumbente. El porcentaje junto a F es el GAP final. Las brechas positivas menores de 0,001\,\% se imprimen como $<0,001$ en lugar de redondearlas a cero. Un guion indica que no existe un GAP comparable. La columna heurística indica unidades pendientes; cero significa demanda completa. Las cargas de planes parciales no se utilizan para calcular mejoras porcentuales.
{\footnotesize\setlength{\tabcolsep}{5pt}
\begin{longtable}{lrrrr}\caption{Resultado de cada problema y GAP final en porcentaje}\label{tab:salidas100}\\
\toprule
ID & Heur.: pendiente & CBC & HiGHS & Gurobi \\
\midrule\endfirsthead
\multicolumn{5}{l}{\tablename\ \thetable\ (continuación)}\\\toprule
ID & Heur.: pendiente & CBC & HiGHS & Gurobi \\
\midrule\endhead\bottomrule\endfoot
E001 & 0 & Ópt. & Ópt. & Ópt. \\
E002 & 0 & Ópt. & Ópt. & Ópt. \\
E003 & 0 & Ópt. & Ópt. & Ópt. \\
E004 & 0 & Ópt. & Ópt. & Ópt. \\
E005 & 0 & Ópt. & Ópt. & Ópt. \\
E006 & 0 & F (0,129) & Ópt. & Ópt. \\
E007 & 0 & Ópt. & Ópt. & Ópt. \\
E008 & 0 & F ($<0,001$) & Ópt. & Ópt. \\
E009 & 0 & Ópt. & Ópt. & Ópt. \\
E010 & 0 & Ópt. & Ópt. & Ópt. \\
E011 & 0 & Ópt. & Ópt. & Ópt. \\
E012 & 0 & Ópt. & Ópt. & Ópt. \\
E013 & 0 & SI & F (0,970) & F (0,015) \\
E014 & 0 & F (1,855) & F (0,067) & Ópt. \\
E015 & 0 & F (9,628) & F (0,385) & F (0,005) \\
E016 & 0 & SI & F (0,580) & F ($<0,001$) \\
E017 & 5788 & Inv. & Inv. & Inv. \\
E018 & 5635 & Inv. & Inv. & Inv. \\
E019 & 7582 & Inv. & Inv. & Inv. \\
E020 & 7269 & Inv. & Inv. & Inv. \\
E021 & 11198 & Inv. & Inv. & Inv. \\
E022 & 13004 & Inv. & Inv. & Inv. \\
E023 & 13420 & Inv. & Inv. & Inv. \\
E024 & 15605 & Inv. & Inv. & Inv. \\
E025 & 0 & F (0,026) & F (0,014) & Ópt. \\
E026 & 0 & F (1,383) & Ópt. & Ópt. \\
E027 & 0 & F (0,380) & F (0,017) & Ópt. \\
E028 & 0 & F (9,067) & F (0,007) & Ópt. \\
E029 & 0 & F (4,240) & F (0,417) & F (0,003) \\
E030 & 0 & SI & F (0,632) & F (0,004) \\
E031 & 0 & F (6,291) & F (0,206) & Ópt. \\
E032 & 0 & SI & F (0,013) & Ópt. \\
E033 & 154 & Ópt. & Ópt. & Ópt. \\
E034 & 10 & F (1,221) & Ópt. & Ópt. \\
E035 & 208 & Ópt. & Ópt. & Ópt. \\
E036 & 132 & F (1,630) & Ópt. & Ópt. \\
E037 & 0 & F (0,778) & F (0,165) & Ópt. \\
E038 & 0 & F (9,310) & F (0,411) & F (0,010) \\
E039 & 0 & SI & F (0,347) & F ($<0,001$) \\
E040 & 0 & F (8,642) & F (0,618) & F (0,006) \\
E041 & 812 & Inv. & Inv. & Inv. \\
E042 & 644 & Inv. & Inv. & Inv. \\
E043 & 896 & Inv. & Inv. & Inv. \\
E044 & 773 & Inv. & Inv. & Inv. \\
E045 & 0 & SI & F (3,541) & F (0,110) \\
E046 & 0 & SI & F (3,026) & F (0,066) \\
E047 & 0 & SI & F (1,187) & F (0,165) \\
E048 & 0 & SI & F (3,523) & F (0,178) \\
E049 & 0 & Ópt. & Ópt. & Ópt. \\
E050 & 0 & Ópt. & Ópt. & Ópt. \\
E051 & 0 & Ópt. & Ópt. & Ópt. \\
E052 & 0 & Ópt. & Ópt. & Ópt. \\
E053 & 0 & F ($<0,001$) & Ópt. & Ópt. \\
E054 & 0 & F (0,209) & F (0,003) & Ópt. \\
E055 & 0 & F ($<0,001$) & Ópt. & Ópt. \\
E056 & 0 & F (1,145) & Ópt. & Ópt. \\
E057 & 0 & F (0,020) & Ópt. & F ($<0,001$) \\
E058 & 0 & F (1,354) & Ópt. & Ópt. \\
E059 & 0 & F (0,150) & Ópt. & Ópt. \\
E060 & 0 & F (0,765) & F (0,333) & Ópt. \\
E061 & 0 & F (4,554) & F (0,415) & F (0,002) \\
E062 & 0 & SI & F (0,874) & F (0,010) \\
E063 & 0 & SI & F (0,938) & F (0,062) \\
E064 & 0 & SI & F (0,749) & F (0,008) \\
E065 & 8870 & Inv. & Inv. & Inv. \\
E066 & 8764 & Inv. & Inv. & Inv. \\
E067 & 9359 & Inv. & Inv. & Inv. \\
E068 & 9236 & Inv. & Inv. & Inv. \\
E069 & 729 & SI & F (0,425) & F ($<0,001$) \\
E070 & 0 & F (3,889) & F (1,209) & F (0,150) \\
E071 & 1232 & SI & F (1,217) & F ($<0,001$) \\
E072 & 335 & F (6,661) & F (1,934) & F (0,060) \\
E073 & 0 & Ópt. & Ópt. & Ópt. \\
E074 & 0 & Ópt. & Ópt. & Ópt. \\
E075 & 0 & Ópt. & Ópt. & Ópt. \\
E076 & 0 & Ópt. & Ópt. & Ópt. \\
E077 & 0 & F (1,455) & F (0,266) & Ópt. \\
E078 & 0 & F (1,314) & F (0,650) & F (0,003) \\
E079 & 0 & F (2,069) & F (0,292) & Ópt. \\
E080 & 0 & F (10,902) & F (0,363) & F (0,009) \\
E081 & 0 & F (3,981) & F (0,464) & F (0,001) \\
E082 & 0 & SI & F (1,722) & F (0,036) \\
E083 & 0 & F (6,180) & F (0,779) & F (0,028) \\
E084 & 460 & F (2,162) & F (1,233) & F (0,034) \\
E085 & 0 & SI & F (2,161) & F (0,050) \\
E086 & 0 & SI & F (2,348) & F (0,068) \\
E087 & 0 & SI & F (3,039) & F (0,179) \\
E088 & 0 & F (13,362) & F (3,170) & F (0,172) \\
E089 & 1364 & Inv. & Inv. & Inv. \\
E090 & 1202 & Inv. & Inv. & Inv. \\
E091 & 1514 & Inv. & Inv. & Inv. \\
E092 & 1409 & Inv. & Inv. & Inv. \\
E093 & 23941 & Inv. & Inv. & Inv. \\
E094 & 24281 & Inv. & Inv. & Inv. \\
E095 & 27374 & Inv. & Inv. & Inv. \\
E096 & 26832 & Inv. & Inv. & Inv. \\
E097 & 0 & SI & F (0,232) & F ($<0,001$) \\
E098 & 0 & SI & F (2,036) & F (0,134) \\
E099 & 0 & SI & F (1,422) & F (0,005) \\
E100 & 0 & SI & F (3,606) & F (0,155) \\
\end{longtable}

}

\subsection{Protocolo de resolución y control de procedencia}\label{an:protocolo100}
\begin{table}[H]\centering\small
\caption{Configuración del estudio de cien instancias}\label{tab:protocolo100}
\begin{tabular}{p{48mm}p{103mm}}\toprule
Elemento & Valor conservado\\\midrule
Generación & Semilla 191007; 25 celdas y cuatro variantes por celda.\\
Remuestreo & Donantes con reemplazo dentro del tercil de venta semestral; supuesto de intercambiabilidad transversal.\\
Ejecuciones & Una por instancia y método; 400 registros.\\
Inicio heredado & No; \texttt{warm\_start=false}.\\
Presupuesto por escala & Estratificado por tamaño y horizonte; igual para los tres motores de cada problema. Los valores exactos figuran en el CSV de resultados.\\
Hilos y semilla de motores & Dos hilos; semilla 42; orden alternado.\\
Modelo & MILP completo; refuerzos válidos comunes; hash del LP idéntico entre motores por instancia.\\
Brechas objetivo & Relativa y absoluta solicitadas iguales a cero.\\
Criterio de certificado & Estado nativo compatible, plan válido, diferencia de cotas $\leq10^{-7}$ min y brecha relativa $\leq10^{-9}$.\\
Plan válido & Residuos de demanda, capacidad, enlace, integridad y dominio $\leq10^{-5}$.\\
Entorno registrado & Versiones, bibliotecas y trabajador registrados en los archivos de la ejecución.\\
CBC & \texttt{feasibilityPump off}, fijado antes del estudio para evitar el fallo nativo identificado en la ejecución previa.\\
Intervalos & Bootstrap percentil por celdas comparables; 1.000 remuestreos; semilla 19042.\\
Lectura del GAP & Brecha relativa entre las cotas finales del mismo problema; se conserva positiva aun cuando sea menor que el redondeo de la tabla.\\\bottomrule
\end{tabular}
\end{table}

La fase principal de la ejecución \href{https://github.com/sebastianporras-web/planeador-milp-rpm/actions/runs/42441274388}{42441274388} produjo los 400 registros de la comparación central. El consolidado, las huellas de las cien entradas y los planes de la heurística acompañan el proyecto editable. Este protocolo estratificado y sus resultados se identifican por esa ejecución y no se combinan con la ejecución anterior de presupuesto uniforme publicada en la rama principal. La revisión cotejó la identidad de las entradas con el diseño original y confirmó 274 planes completos entre los cuatro métodos y 92 cierres numéricos entre los tres motores. El flujo global registró un fallo posterior al exportar cuatro diagnósticos auxiliares de inviabilidad; esa salida no forma parte de los 400 registros primarios ni de las tablas de este estudio.

El caso E097 conserva un GAP positivo pese a la cercanía de sus cotas: la diferencia absoluta excede el umbral conservador de la auditoría. La tabla principal lo presenta como plan completo con brecha residual, mientras el inventario mantiene su identidad de entrada.
